\section{Proofs for Multiplicity and Sequentially Rejective Testing}
\label{annex:p15-multiple-testing}

\begin{annexitem}{annex:p15.mt.inclusive-levels}{thm:p15-mt-inclusive-levels}{Validity of every inclusive supremum tail}
\begin{proof}
Fix \(i\) and \(\theta\in\omega_i\). Write
\(G(y)=P_\theta(Y_i\geq y)\). It is nonincreasing and Borel measurable.
For \(0\leq u<1\), the set \(A_u=\{y:G(y)\leq u\}\) is an upper ray
or empty. If empty, \(P_\theta\{G(Y_i)\leq u\}=0\).
If nonempty, its infimum \(a\) is finite: \(G(y)\to1\) as
\(y\to-\infty\), so \(A_u\) is bounded below.
If \(a\in A_u\), the event is \(\{Y_i\geq a\}\) and its probability
is \(G(a)\leq u\). Otherwise it is \(\{Y_i>a\}\).
Choose \(a_j\in A_u\) decreasing to \(a\); continuity from below gives
\[
 P_\theta(Y_i>a)=\lim_jP_\theta(Y_i\geq a_j)\leq u.
\]
This includes \(u=0\). At \(u=1\) the inequality is immediate.
Thus \(G(Y_i)\) is superuniform for every real-valued law.

By the definition of the supremum, \(\alpha_i(y)\geq G(y)\).
Consequently \(\{\alpha_i(Y_i)\leq u\}\subseteq\{G(Y_i)\leq u\}\),
which proves (V). The supremum of nonincreasing functions is again
nonincreasing, hence Borel; composing with the measurable statistic
gives a measurable level. The argument used just the particular actual
null law among the supremum's candidates; neither atomlessness nor a
joint relation between different tests entered.
\end{proof}
\end{annexitem}

\begin{annexitem}{annex:p15.mt.ordering}{thm:p15-mt-ordering}{Every sorting, equality and stopping boundary}
\begin{proof}
For each of the finitely many permutations, the event that it is the
chosen ordering is a finite intersection of comparisons \(R_i<R_j\)
or \(R_i\leq R_j\), chosen according to the fixed index tie rule.
These events are measurable and partition the sample space.
On each such event, rejection of its \(r\)-th hypothesis is the finite
intersection of \(R^{(j)}\leq\alpha/(n-j+1)\), \(1\leq j\leq r\).
A finite union over ordering events proves measurability of each
original index's rejection event.

The index either increases or the algorithm stops, so at most \(n\)
comparisons occur. If a tied value \(v\) occupies positions \(a,\ldots,b\),
and the procedure reaches \(a\), failure at \(a\) stops before the entire
block. If \(v\leq\alpha/(n-a+1)\), every later threshold within the
block is at least this one, so every member passes. Permuting that block
changes neither its starting position nor the hypotheses outside it.
If the procedure had stopped earlier, none of its members is rejected
under any tie order. Equality is included by the displayed nonstrict
comparison. For \(n=1\) its denominator is one; the explicit empty-family
convention does not make a division or a comparison.
\end{proof}
\end{annexitem}

\begin{annexitem}{annex:p15.mt.holm-strong}{thm:p15-mt-holm-strong}{The first true rejection under an arbitrary configuration}
\begin{proof}
Fix the full parameter \(\theta\), including all nuisance components.
If \(I_0(\theta)\) is empty the error probability is zero. Otherwise
let \(m=\abs{I_0(\theta)}\geq1\).
If a true hypothesis is rejected, let \(j\) be the smallest rank of a
true hypothesis. Every hypothesis preceding it is false, so
\(j-1\leq n-m\), or \(n-j+1\geq m\).
The procedure could reject that true hypothesis only after reaching and
passing its comparison. Therefore
\[
 R^{(j)}\leq\frac{\alpha}{n-j+1}\leq\frac{\alpha}{m}.
\]
This level belongs to a true index. Hence the error event is contained
in the union of the events \(R_i\leq\alpha/m\), \(i\in I_0(\theta)\).
The finite union inequality and the actual marginal premises give
\[
 P_\theta\{\mathcal R_{\rm H}\cap I_0(\theta)\ne\varnothing\}
 \leq\sum_{i\in I_0(\theta)}P_\theta(R_i\leq\alpha/m)
 \leq m\frac{\alpha}{m}=\alpha.
\]
No distribution was imposed on a false-hypothesis level. The argument
allows it to depend arbitrarily on any true level, and allows arbitrary
dependence among true levels.

For a specified nonempty \(I\) and parameter in its intersection null,
\(I\subseteq I_0(\theta)\). Rejecting a member of \(I\) is therefore
contained in the error event just bounded. Conversely, the source's
subset statement applied to \(I=I_0(\theta)\) gives strong control
whenever this set is nonempty; the empty case is automatic. This proves
the equivalence of the two protection formulations, whether or not
every abstract truth configuration is attainable in the parameter model.
\end{proof}
\end{annexitem}

\begin{annexitem}{annex:p15.mt.domination}{thm:p15-mt-domination}{The same data, three different rejection sets}
\begin{proof}
The indices satisfying \(R_i\leq\alpha/n\) form an initial segment in
the obtained-level ordering. Every comparison in that segment passes
Holm because \(\alpha/(n-j+1)\geq\alpha/n\). Thus
\(\mathcal R_{\rm B}\subseteq\mathcal R_{\rm H}\).
Every Holm-rejected level is at most its own threshold, and that
threshold is at most \(\alpha\). This gives the other inclusion.
For fixed Bonferroni,
\[
 P_\theta(\mathcal R_{\rm B}\cap I_0\ne\varnothing)
 \leq\sum_{i\in I_0}P_\theta(R_i\leq\alpha/n)
 \leq m\alpha/n\leq\alpha.
\]

Pointwise inclusion implies inclusion of the events \(J\subseteq\mathcal R\)
and \(J\cap\mathcal R\ne\varnothing\), and the pointwise inequality
\(\abs{J\cap\mathcal R_{\rm B}}\leq\abs{J\cap\mathcal R_{\rm H}}\).
Taking probabilities or expectations proves the three power comparisons
under any specified alternative law. If \(J\) is empty the same
conventions give probabilities one and zero and count zero, respectively.
Equality can occur, for example when \(n=1\), or when the levels have
support on outcomes where the two rules coincide.

Holm makes any rejection exactly when its first comparison passes:
\(\min_i R_i\leq\alpha/n\). This is also Bonferroni's any-rejection
event, so their no-rejection events coincide at every law, not just
at the complete null.
\end{proof}
\end{annexitem}

\begin{annexitem}{annex:p15.mt.adjusted}{thm:p15-mt-adjusted}{The running maximum and tied values}
\begin{proof}
The \(r\)-th hypothesis is rejected exactly when
\((n-j+1)R^{(j)}\leq\alpha\) for every \(j\leq r\).
That condition is the maximum inequality in the definition.
Since \(\alpha<1\), capping the maximum at one does not change whether
it is at most \(\alpha\). Taking the maximum over a growing index set
shows nondecrease in \(r\).

If a value \(v\geq0\) fills ranks \(a,\ldots,b\), the products within
that block are \((n-a+1)v,\ldots,(n-b+1)v\), a nonincreasing list.
The largest block contribution already enters at \(a\). The running
maximum, and its cap, therefore stay the same throughout the block.
This proves adjusted-level invariance as well as decision invariance.
For the displayed five-level example the products are
\(0.020,0.048,0.054,0.038,0.040\); their running maxima are the stated
\(0.020,0.048,0.054,0.054,0.054\).
\end{proof}
\end{annexitem}

\begin{annexitem}{annex:p15.mt.intersection-policy}{thm:p15-mt-intersection-policy}{Exactly the specified logical implications}
\begin{proof}
If \(I\subseteq J\) and \(I\cap\mathcal R\ne\varnothing\), then
\(J\cap\mathcal R\ne\varnothing\); this is coherence for the stated
containment implication. The witness \(i\in I\cap\mathcal R\) is a
rejected basic constituent, which proves consonance.
If a secondary \(H_I\) is true, every \(i\in I\) is a true basic index.
Its rejection therefore implies rejection of a true basic member.
Conversely a rejected true basic member is a falsely rejected true
singleton secondary. Thus, over the full nonempty intersection
inventory, the two family error events are identical.

Neither index containment nor this error identity asserts a relation
between semantically equivalent or otherwise dependent index sets.
The explicit restriction in the body leaves those further consistency
conditions separate rather than treating them as consequences of the
set-theoretic argument.
\end{proof}
\end{annexitem}

\begin{annexitem}{annex:p15.mt.closed-bonferroni}{thm:p15-mt-closed-bonferroni}{Every intersection test and the exact basic decision}
\begin{proof}
Under the intersection null for \(I\), each member is marginally valid,
so the local test has rejection probability at most
\(\sum_{i\in I}\alpha/\abs I=\alpha\).
Now fix the observed ordering. Suppose Holm rejects a hypothesis of
rank \(r\), and let \(I\) be any index set containing it. Let \(k\leq r\)
be the smallest rank represented in \(I\). Then
\(\abs I\leq n-k+1\) and
\[
 \min_{i\in I}R_i=R^{(k)}
        \leq\frac{\alpha}{n-k+1}\leq\frac{\alpha}{\abs I}.
\]
That local test rejects. All containing intersections therefore reject.

Conversely, if Holm does not reject the chosen hypothesis, it stopped
at a first failing rank \(k\) no later than that hypothesis's rank.
Take \(I=\{i_k,\ldots,i_n\}\), which contains the chosen index.
Its smallest level is \(R^{(k)}>\alpha/(n-k+1)=\alpha/\abs I\),
so this containing local test does not reject. The closed rule also
does not reject the chosen basic hypothesis.
If Holm rejected all hypotheses the first direction already covers all
indices. The nonstrict passing and strict failing boundaries treat ties
and equality exactly.
\end{proof}
\end{annexitem}

\begin{annexitem}{annex:p15.mt.directional}{thm:p15-mt-directional}{The common boundary law and the two conclusions}
\begin{proof}
Take any one of the stated common boundary laws \(P_*\) at
\(\gamma=\gamma_0\), with its nuisance values fixed. That law belongs
to both null sets. Consequently the two supremum tails satisfy
\[
 R_a(y)+R_b(-y)
  \geq P_*(Y\geq y)+P_*(-Y\geq-y)
  =P_*(Y\geq y)+P_*(Y\leq y)
  =1+P_*(Y=y)\geq1.
\]
This uses inclusive tails; an atom strengthens the inequality.

Suppose both members are rejected by ordinary Holm. When the first
member is tested the other is still in the remaining family, so its
level is at most \(\alpha/2\). The other rejected level is at most
\(\alpha\). Their sum is therefore at most \(3\alpha/2\leq3/4\),
contradicting the tail inequality.

Applying the argument to every specified pair precludes all opposite
pair conclusions. The hypothesis that no other logical constraints
are at issue is what permits the further statement that these
contradictions exhaust the stated consistency problem.
\end{proof}
\end{annexitem}

\begin{annexitem}{annex:p15.mt.confidence-identical}{thm:p15-mt-confidence-identical}{No rejection and inversion at every candidate}
\begin{proof}
For each fixed observation and candidate \(v\), the procedure makes
no rejection exactly when its first comparison fails:
\(\min_iR_i(x;v)>\alpha/n\). This is equivalent to every displayed
strict inequality, and is exactly Bonferroni's no-rejection condition.
Thus the two candidate sets coincide pointwise, including threshold
equalities, which are rejected and therefore excluded.
Their membership graphs are measurable finite intersections because
the candidate-dependent level functions are jointly measurable.

At the actual candidate, all the detail hypotheses are true. By the
strong-control theorem, their any-rejection probability is at most
\(\alpha\). Its complementary event is precisely membership of that
candidate in the inverted set, giving coverage at least \(1-\alpha\).
This reasoning specifies one inversion construction and proves no
identity for other interval or confidence-region procedures.
\end{proof}
\end{annexitem}

\begin{annexitem}{annex:p15.mt.normal-power}{thm:p15-mt-normal-power}{A complete finite probability recursion for the source's example}
\begin{proof}
For \(Y\sim N(\mu,1)\), continuity and monotonicity of \(\Phi\) give
\[
 P(1-\Phi(Y)\leq u)
 =P\{Y\geq\Phi^{-1}(1-u)\}
 =\Phi\bigl(\mu-\Phi^{-1}(1-u)\bigr).
\]
The obtained level is strictly between zero and one almost surely,
which gives the two endpoint values.
Bonferroni rejects each mean-three hypothesis if
\(Y_i\geq z\), \(z=\Phi^{-1}(0.995)\).
Their independence gives \(\Phi(3-z)^2\).

Here is the finite recursion for Holm. Label the three groups by
\(g=0,6,3\), with respective totals
\(M_0=4,M_6=4,M_3=2\). For bin \(k=1,\ldots,10\), define
\[
 q_{g,k}=
 \frac{F_g(b_k)-F_g(b_{k-1})}{1-F_g(b_{k-1})}.
\]
Every denominator is positive. If a state has already placed
\(a_g\) of group \(g\)'s \(M_g\) levels in earlier bins, its conditional
probability of placing \(t_g\) in the current bin is
\[
 K_k(a,t)=\prod_{g\in\{0,6,3\}}
       \binom{M_g-a_g}{t_g}
       q_{g,k}^{t_g}(1-q_{g,k})^{M_g-a_g-t_g},
       \quad0\leq t_g\leq M_g-a_g.                            \tag{K}
\]
The three binomial counts are independent. To see this directly,
condition on which indices appeared in previous bins. Each remaining
independent level has its group's law conditioned on being greater
than \(b_{k-1}\); hence its current-bin probability is \(q_{g,k}\).
Summing over the possible previous index identities gives the same
binomial count probabilities, since all members of a group have the
same law. Previous success depends only on earlier bin counts and
does not further constrain these remaining values.

Let \(D_k(a)\) be the probability mass after bins \(1,\ldots,k\) of
reaching the placed-count vector \(a=(a_0,a_6,a_3)\), with no stopping
failure and with \(a_3<2\). Start with \(D_0(0,0,0)=1\) and all other
masses zero. Keep a cumulative success mass \(B_0=0\).
For each state in \(D_{k-1}\) and each permitted \(t\), form the
transition mass \(D_{k-1}(a)K_k(a,t)\).
If \(\sum_gt_g>0\) and \(\sum_ga_g+1<k\), discard this mass as a
stopping failure. Otherwise let \(a'=a+t\). If \(a'_3=2\), add the
mass to \(B_k\), carrying forward the existing \(B_{k-1}\); otherwise
add it to \(D_k(a')\). Initialize each new \(D_k\) to zero before
adding its transitions.

Indeed, the first rank in a currently occupied bin is
\(r=\sum_ga_g+1\). If \(r<k\), its level exceeds \(b_{k-1}\geq b_r\),
so it fails the rank-\(r\) comparison. If \(r\geq k\), its level is
at most \(b_k\leq b_r\), so it passes, and every subsequent member
of that bin has a still larger rank and threshold. Empty bins change
no rank and cause no comparison. This proves the complete transition
rule, not just a necessary condition for rejection.
Absorbed success remains success regardless of what later tests do.
At \(k=10\), levels not yet placed exceed \(b_{10}=\alpha\);
no remaining mean-three level can ever be rejected. Thus the required
probability is exactly \(B_{10}\), a finite sum determined by (K).

Evaluation with \(z\simeq2.5758293035\) gives
\(\Phi(3-z)^2\simeq0.4412670349\).
Evaluating the stated finite recurrence gives
\(B_{10}\simeq0.5631999191\). The formula and bin construction are
exact; the final decimal evaluations use ordinary numerical normal
CDF and quantile values. They are not a distribution-free power claim.
\end{proof}
\end{annexitem}

\begin{annexitem}{annex:p15.mt.weighted}{thm:p15-mt-weighted}{Every true configuration with positive fixed constants}
\begin{proof}
Fix \(\theta\). With no true indices the FWER is zero.
Otherwise put \(C_0=\sum_{i\in I_0(\theta)}c_i>0\).
If any true hypothesis is rejected, inspect the first such member in
the ratio ordering. All true hypotheses, including it, are still
present in that comparison's remaining weight, so \(W_j\geq C_0\).
Passing gives
\[
 S_{i_j}\leq\alpha/W_j\leq\alpha/C_0,\qquad
 R_{i_j}\leq\alpha c_{i_j}/C_0.
\]
Every allocation on the right is between zero and \(\alpha<1\).
The marginal bounds and union inequality therefore give
\[
 \operatorname{FWER}_\theta(\mathcal R_{\rm W})
 \leq\sum_{i\in I_0}P_\theta(R_i\leq\alpha c_i/C_0)
 \leq\sum_{i\in I_0}\alpha c_i/C_0=\alpha.
\]
The constants must be fixed for this argument; the allocation thresholds
are not chosen by examining the true levels.

Multiplying all \(c_i\) by \(\lambda>0\) divides every \(S_i\) by
\(\lambda\) and multiplies every \(W_j\) by \(\lambda\). Both the
ordering and comparison are unchanged. Equal constants \(c\) yield
\(R_{i_j}/c\leq\alpha/((n-j+1)c)\), precisely ordinary Holm.

For a tied ratio block, its starting remaining total is independent
of the ordering within the block. If the first ratio fails, the whole
block remains unrejected. If it passes, removing any positive block
weight decreases \(W_j\), so the threshold increases and all other
equal ratios pass. Removing the whole block leaves the same weight
and indices for the rest of the algorithm. This proves tie invariance.
Finite sorting and comparisons of measurable ratios also prove the
weighted procedure's measurability and finite termination.

The fixed weighted Bonferroni set is the initial segment with
\(S_i\leq\alpha/\sum_\ell c_\ell\). At every step \(W_j\) is at most
that original sum, so every member of the segment passes.
It is therefore contained in the weighted step-down rejection set.
For the first three-level example, the ratios are
\(0.00375,0.04,0.04\); the first threshold is \(0.05/10=0.005\),
and the second is \(0.05/2=0.025\), giving precisely the stated
single rejection. For the second example the first ratio is
\(0.048/8=0.006>0.005\), so it stops at once; ordinary Holm's raw
sorted levels \(0.008,0.009,0.048\) pass all three thresholds.
These computations rule out a claimed universal dominance over
unweighted Holm.

Under the opposite-tail conditions previously proved, the two raw
levels have sum at least one. Every weighted-rejected raw level is
at most \(\alpha\), since its weight is at most the remaining total.
For whichever of the pair is tested first, the other positive weight
is still present, so this first raw level is strictly below \(\alpha\).
Rejecting both at \(\alpha\leq1/2\) would therefore give a sum strictly
below \(2\alpha\leq1\), a contradiction.
\end{proof}
\end{annexitem}

\begin{annexitem}{annex:p15.mt.independent}{thm:p15-mt-independent}{True-null independence and each threshold comparison}
\begin{proof}
Let \(m\geq1\) true hypotheses be present at the fixed parameter.
At the first true rank, at least \(m\) hypotheses remain. For
\(0<1-\alpha<1\), \(d_k=1-(1-\alpha)^{1/k}\) decreases as \(k\)
increases. A true rejection therefore implies that some true level
is at most \(d_m\). Mutual independence of these true levels gives
\[
 P_\theta\{\min_{i\in I_0}R_i\leq d_m\}
 =1-\prod_{i\in I_0}P_\theta(R_i>d_m)
 \leq1-(1-d_m)^m=\alpha.
\]
The all-false case has error probability zero. No independence of
false levels was used.

For \(k=1\), \(d_1=\alpha\). For \(k>1\) put \(x=\alpha/k\).
The function \(f(x)=(1-x)^k-1+kx\) has \(f(0)=0\) and
derivative \(k[1-(1-x)^{k-1}]>0\) for \(0<x<1\).
Thus \((1-\alpha/k)^k>1-\alpha\), and taking positive \(k\)-th
roots gives \(d_k>\alpha/k\).
Every ordinary Holm-passing prefix consequently also passes this
independent-tail rule with the same raw ordering and stopping convention.
The dependence counterexample below shows that the independent product
calculation is an essential premise for this larger-threshold claim.
\end{proof}
\end{annexitem}

\begin{annexitem}{annex:p15.mt.control-model}{thm:p15-mt-control-model}{Pooled normal coordinates and every one-sided null}
\begin{proof}
Apply the orthogonal-coordinate construction of
\cref{thm:p15-student-orthogonal} to each group with \(N_g\geq2\).
For \(N_g=1\), the single standardized observation is already standard
normal and its residual coordinate list is empty.
Its standardized group mean is standard normal and is independent of
its \(N_g-1\) standardized residual coordinates.
Group independence makes all these orthogonal coordinates independent.
Therefore
\[
 W=\nu s_{\rm p}^2/\sigma^2
\]
is a sum of \(\nu\) independent standard-normal squares and has
\(\chi^2_\nu\) law, independently of every group mean.
It is strictly positive almost surely because the specified \(\nu\) is positive.
For the fixed comparison \(i\),
\[
 Z_i=\frac{(\bar X_i-\mu_i)-(\bar X_0-\mu_0)}
                  {\sigma\sqrt{1/N_i+1/N_0}}
\]
is standard normal, independently of \(W\), by the independent-normal
sum law. Put
\(\lambda_i=(\mu_i-\mu_0)/(\sigma\sqrt{1/N_i+1/N_0})\).
Then \(T_i=(Z_i+\lambda_i)/\sqrt{W/\nu}\) almost surely.
Under \(H_i\), \(\lambda_i\leq0\), so with the same \(Z_i,W\) it is
pointwise at most \(Z_i/\sqrt{W/\nu}\), whose continuous Student CDF
is \(F_\nu\) by \cref{thm:p15-student-law}.
For \(0<u<1\),
\[
 P_\theta\{1-F_\nu(T_i)\leq u\}
 =P_\theta\{T_i\geq t_{\nu,1-u}\}\leq u.
\]
The bound at \(u=0\) follows from finiteness of the statistic almost
surely, and the bound at \(u=1\) is trivial. Defining \(R_i=1\) on
the zero-scale set changes none of these probabilities and preserves
measurability. This establishes (V) for every comparison, including all
other means and the nuisance variance.

The different comparisons need not be independent: their standardized
numerators share the control mean and all denominators share \(W\).
The already proved union-bound procedures nevertheless apply.
For ten groups of four observations, \(\nu=40-10=30\).
At \(k\) remaining comparisons the upper-tail condition
\(R_i\leq0.05/k\) is exactly \(T_i\geq t_{30,1-0.05/k}\), giving
the stated order of the comparisons. This derives the Holm side
of the source's table; it makes no derivation claim for its separately
reported refined Dunnett entries.
\end{proof}
\end{annexitem}

\begin{annexitem}{annex:p15.mt.conditional-calibration}{thm:p15-mt-conditional-calibration}{A finite conditional law and integration}
\begin{proof}
Taking expectations of the assumed conditional inequality yields
\[
 P_\theta(R_i\leq u)
   =E_\theta[P_\theta(R_i\leq u\mid Z_i)]\leq u.
\]
For the finite permutation case, condition on the stated pooled
information. Suppose the allowed assignments form a finite set
\(\mathcal A\) of size \(M\geq1\), each with conditional probability
\(1/M\). For the ordered statistic \(T\), give assignment \(a\)
the inclusive level
\[
       r(a)=\frac{\abs{\{b\in\mathcal A:T(b)\geq T(a)\}}}{M}.
\]
List the distinct values of \(T\) from largest downward, with their
positive integer multiplicities \(m_1,m_2,\ldots\).
Every assignment in block \(j\) has level
\((m_1+\cdots+m_j)/M\). The blocks whose level is at most \(u\)
form an initial segment, whose total multiplicity is at most \(Mu\).
If there is no such block the count is zero. Thus the conditional
rejection probability is at most \(u\), even with tied statistics.
Integrating proves marginal validity.

Uniformity over assignments is the stated null exchangeability premise,
not a consequence of observing a finite dataset. Different tests may
use different finite conditional laws and different pooled information.
Their separate marginal inequalities suffice for the previously proved
family control; the conditional argument adds no independence premise.
\end{proof}
\end{annexitem}

\begin{annexitem}{annex:p15.mt.failures}{thm:p15-mt-failures}{The exact probabilities of four failures and their remedies}
\begin{proof}
For independent uniform true levels, the probability that all \(n\)
levels exceed \(t\) is \((1-t)^n\). Applying \(t=0.05\) to \(n=20\)
gives \(1-0.95^{20}\). For both Holm and Bonferroni, any rejection
means the minimum is at most \(0.05/20\), giving the second formula.
Selecting the smallest level and testing it at the single-test cutoff
is the first event again. The selected minimum is not a valid marginal
single-test p-value: its CDF is \(1-(1-u)^{20}>u\) for \(0<u<1\).
Retaining the actual twenty-member family restores the proved control.

For the second construction make the experiment explicit.
Use parameter space \(\{\theta_0,\theta_1\}\);
let \(H_1\) have null set \(\{\theta_0\}\) and \(H_2\) have null
set \(\{\theta_0,\theta_1\}\). Under \(\theta_0\), \(R_1,R_2\)
are independent uniforms; under \(\theta_1\), take \(R_1=0\)
and \(R_2\) uniform. Every level is valid under every law where its
own hypothesis is true. The whole-family gate has complete-null error
\[
           1-(1-\alpha/2)^2=\alpha-\alpha^2/4\leq\alpha.
\]
Under \(\theta_1\) it always opens and falsely rejects the true \(H_2\),
so its FWER is one. Ordinary Holm instead tests \(H_2\), after the
certain rejection of \(H_1\), at level \(\alpha\), so its error is
\(\alpha\) there. This checks the precise configuration that the
weaker gate guarantee missed.

For \(R_1=U,R_2=1-U\), with \(U\) uniform, put
\(d=1-\sqrt{1-\alpha}\). At \(\alpha=0.05\), \(0<d<1/2\).
The two intervals \(U\leq d\) and \(U\geq1-d\) are disjoint apart
from null endpoints and have total probability \(2d\).
They are precisely the independent-tail procedure's any-rejection
event. The value \(2(1-\sqrt{0.95})>0.05\) establishes the failure.
Ordinary Holm uses \(d=\alpha/2\) at its first step, yielding total
probability \(2(\alpha/2)=\alpha\).

Finally choose the two positive weights after seeing the smaller level.
It always has the smaller ratio: for \(u\leq v\),
\(u/99\leq v\). It is tested first at raw threshold
\(99\alpha/100\). Thus any rejection is exactly
\(\min(R_1,R_2)\leq99\alpha/100\), and independent uniforms give
\[
 1-(1-99\alpha/100)^2.
\]
At \(\alpha=0.05\) this is \(0.09654975>0.05\).
The error comes from correlating the allocation with the observed
levels. Fixing the weights and the family before these observations
returns to the weighted theorem's hypotheses.
\end{proof}
\end{annexitem}

\begin{annexitem}{annex:p15.mt.calibration-error}{thm:p15-mt-calibration-error}{Every marginal error contribution}
\begin{proof}
With no true hypotheses the error probability is zero.
Otherwise the ordinary first-true argument gives
\[
 \operatorname{FWER}_\theta
 \leq\sum_{i\in I_0}P_\theta(R_i\leq\alpha/m)
 \leq\sum_{i\in I_0}(\alpha/m+\varepsilon_i)
 =\alpha+\sum_{i\in I_0}\varepsilon_i.
\]
For fixed positive weights, let \(C_0=\sum_{i\in I_0}c_i\).
The weighted first-true argument replaces the allocation \(\alpha/m\)
by \(\alpha c_i/C_0\) for index \(i\). These allocations again sum
to \(\alpha\), so the same bound results.
Every probability is also at most one; taking the minimum yields
the displayed statement. This explicitly retains the calibration
residual instead of asserting the exact guarantee for approximate inputs.
\end{proof}
\end{annexitem}
